\documentclass[11pt,letterpaper]{article}

\usepackage[utf8]{inputenc}
\usepackage[T1]{fontenc}
\usepackage[margin=0.5in]{geometry}
\usepackage{lmodern}
\usepackage{microtype}
\usepackage{amsmath}
\usepackage{amssymb}
\usepackage{mathtools}
\usepackage{enumitem}
\usepackage[hidelinks]{hyperref}

\setlength{\parindent}{0pt}
\setlength{\parskip}{0.45em}
\setlength{\abovedisplayskip}{0.7em}
\setlength{\belowdisplayskip}{0.7em}
\setlist[itemize]{leftmargin=1.5em, itemsep=0.15em, topsep=0.2em}
\setlist[enumerate]{leftmargin=1.5em, itemsep=0.15em, topsep=0.2em}

\begin{document}

\begin{center}
{\Large\bfseries Unpacking the Camera-to-Screen Pipeline and Jacobian in 3D Gaussian Splatting\par}
\vspace{0.4em}
\end{center}

\textbf{Let's unpack every single step} from a \textbf{camera-space 3D point} to a \textbf{2D screen point}, and then show \textbf{exactly where the Jacobian enters}.

I'll do it in the cleanest order:

\begin{enumerate}
\item start with a 3D point in camera coordinates
\item send it through the graphics pipeline step by step
\item define the full map $(h)$
\item explain the local linearization
\item show how covariance uses that
\end{enumerate}

\section{What We Are Trying to Compute}

Suppose you already have a Gaussian center in \textbf{camera coordinates}:

\[
x_c =
\begin{bmatrix}
X_c\\
Y_c\\
Z_c
\end{bmatrix}
\]

This is a 3D point relative to the camera.

We want to know:

\begin{itemize}
\item where this point lands on the screen
\item how a \textbf{small cloud of nearby points} around it gets turned into a 2D ellipse
\end{itemize}

That second part is where covariance comes in.

\section{Camera Coordinates}

At this stage, the point is in the camera's own coordinate system.

So:

\begin{itemize}
\item $(X_c)$: left-right location relative to camera
\item $(Y_c)$: up-down location relative to camera
\item $(Z_c)$: depth or distance in front of the camera
\end{itemize}

This point is still 3D.

\section{Convert to Homogeneous Coordinates}

In graphics, points are often written in homogeneous form:

\[
\tilde{x}_c =
\begin{bmatrix}
X_c\\
Y_c\\
Z_c\\
1
\end{bmatrix}
\]

Why? Because matrix multiplication with $(4 \times 4)$ matrices becomes convenient.

\section{Apply the Perspective Projection Matrix}

Now we multiply by the projection matrix $(P)$:

\[
\tilde{x}_{\text{clip}} = P \tilde{x}_c
\]

This gives the \textbf{clip-space} point:

\[
\tilde{x}_{\text{clip}} =
\begin{bmatrix}
x_{\text{clip}}\\
y_{\text{clip}}\\
z_{\text{clip}}\\
w_{\text{clip}}
\end{bmatrix}
\]

This is still \textbf{not} the final screen coordinate. It is an intermediate representation.

\subsection{What Does the Projection Matrix Look Like?}

For a symmetric perspective camera, one common form is:

\[
P =
\begin{bmatrix}
\frac{1}{\tan(\theta_x/2)} & 0 & 0 & 0\\
0 & \frac{1}{\tan(\theta_y/2)} & 0 & 0\\
0 & 0 & a & b\\
0 & 0 & -1 & 0
\end{bmatrix}
\]

Or equivalently, in pixel-style intrinsics language, people often think in terms of:

\begin{itemize}
\item focal lengths
\item principal point
\item perspective divide
\end{itemize}

The exact matrix form depends on convention, but the important point is \textbf{this step prepares the point for the perspective divide}.

\section{Perspective Divide}

This is the crucial step.

After projection matrix multiplication, you divide by $(w_{\text{clip}})$:

\[
x_{\text{ndc}} = \frac{x_{\text{clip}}}{w_{\text{clip}}},\qquad
y_{\text{ndc}} = \frac{y_{\text{clip}}}{w_{\text{clip}}},\qquad
z_{\text{ndc}} = \frac{z_{\text{clip}}}{w_{\text{clip}}}
\]

These are the coordinates in \textbf{NDC} (normalized device coordinates).

This is where the mapping becomes nonlinear, because of the division.

That is the key step you were asking about earlier.

\section{NDC Coordinates}

Now the point is inside a normalized cube:

\begin{itemize}
\item typically $(x_{\text{ndc}} \in [-1,1])$
\item $(y_{\text{ndc}} \in [-1,1])$
\end{itemize}

This is often called the \textbf{canonical view volume} or \textbf{NDC space}.

So after the perspective divide, the point is no longer in real camera units. It is in normalized coordinates.

\section{Viewport Transform: \texorpdfstring{NDC $\to$ Screen Pixels}{NDC to Screen Pixels}}

Now convert NDC to actual screen pixels.

If the image width is $(W)$ and the height is $(H)$, then one common mapping is:

\[
u = \frac{W}{2}(x_{\text{ndc}} + 1)
\]

\[
v = \frac{H}{2}(1 - y_{\text{ndc}})
\]

Now $(u,v)$ is the actual screen pixel position.

So the pipeline so far is:

\[
x_c \to \tilde{x}_c \to \tilde{x}_{\text{clip}} \to x_{\text{ndc}} \to s
\]

where

\[
s =
\begin{bmatrix}
u\\
v
\end{bmatrix}
\]

\section{The Full Map \texorpdfstring{$(h)$}{(h)}}

Now let's bundle all those steps together.

Define one function:

\[
s = h(x_c)
\]

That means:

\[
h : \mathbb{R}^3 \to \mathbb{R}^2
\]

and it takes a 3D camera-space point $(x_c)$ and returns the 2D screen point $(s)$.

So when I said earlier

\[
s = h(x_c)
\]

I meant \textbf{all the little graphics steps together, treated as one big function}.

That function $(h)$ includes:

\begin{enumerate}
\item homogeneous conversion
\item projection matrix
\item perspective divide
\item NDC to screen conversion
\end{enumerate}

\section{Why \texorpdfstring{$(h)$}{(h)} Is Not Linear}

Because inside $(h)$, there is a division by $(w_{\text{clip}})$.

That is the same reason perspective projection was nonlinear before.

So $(h)$ is not a simple matrix multiplication from $(\mathbb{R}^3 \to \mathbb{R}^2)$.

That is why we cannot just write:

\[
\Sigma_{\text{screen}} = A \Sigma_c A^T
\]

with one fixed $(A)$ for all points.

\section{Small Offsets Around the Gaussian Center}

Now suppose the Gaussian center is at

\[
\mu_c
\]

A nearby point in the Gaussian is:

\[
\mu_c + \delta x
\]

where $(\delta x)$ is a tiny 3D offset.

This means we are no longer looking at just one point. We are looking at a tiny neighborhood around the center.

\section{Exact Screen Mapping of That Nearby Point}

The exact screen location of that nearby point is:

\[
h(\mu_c + \delta x)
\]

But this exact function is nonlinear and messy.

So we approximate it.

\section{Local Linear Approximation}

Near the center $(\mu_c)$, the function $(h)$ behaves approximately like a linear map.

So:

\[
h(\mu_c + \delta x) \approx h(\mu_c) + J_h \, \delta x
\]

This is the important formula.

Let's decode every part.

\subsection{\texorpdfstring{$h(\mu_c)$}{h(mu\_c)}}

This is the screen position of the Gaussian center.

\subsection{\texorpdfstring{$\delta x$}{delta x}}

Tiny 3D offset around the center.

\subsection{\texorpdfstring{$J_h$}{J\_h}}

The Jacobian of the full camera-to-screen map $(h)$, evaluated at $(\mu_c)$.

It is a $(2 \times 3)$ matrix.

\subsection{\texorpdfstring{$J_h \delta x$}{J\_h delta x}}

Approximate 2D screen offset caused by the tiny 3D offset.

So:

\[
\delta s \approx J_h \delta x
\]

where

\[
\delta s = h(\mu_c + \delta x) - h(\mu_c)
\]

\section{What Is the Jacobian \texorpdfstring{$(J_h)$}{(J\_h)}?}

It is the matrix of partial derivatives:

\[
J_h =
\begin{bmatrix}
\frac{\partial u}{\partial X_c} & \frac{\partial u}{\partial Y_c} & \frac{\partial u}{\partial Z_c}\\
\frac{\partial v}{\partial X_c} & \frac{\partial v}{\partial Y_c} & \frac{\partial v}{\partial Z_c}
\end{bmatrix}
\]

This tells you:

\begin{itemize}
\item how screen $(u)$ changes if $(X_c)$ changes a tiny bit
\item how screen $(u)$ changes if $(Y_c)$ changes a tiny bit
\item how screen $(u)$ changes if $(Z_c)$ changes a tiny bit
\item the same for $(v)$
\end{itemize}

So $(J_h)$ is a \textbf{tiny-motion converter}:

\[
\text{tiny 3D camera offset} \to \text{tiny 2D screen offset}
\]

\section{Why Covariance Now Uses the Standard Rule}

Once we write

\[
\delta s \approx J_h \delta x
\]

the nearby offsets are now approximately related by a \textbf{linear map}.

That is the key.

For a linear map

\[
y = A x
\]

covariance transforms as

\[
\Sigma_y = A \Sigma_x A^T
\]

Here:

\begin{itemize}
\item $(x = \delta x)$
\item $(y = \delta s)$
\item $(A = J_h)$
\end{itemize}

So:

\[
\Sigma_{\text{screen}} \approx J_h \Sigma_c J_h^T
\]

That is the covariance formula.

\section{So Where Exactly Are We in the Pipeline When We Do This?}

We are at this stage:

\textbf{First:} world Gaussian $\to$ camera Gaussian

That gives $(\Sigma_c)$.

\textbf{Then:} camera Gaussian $\to$ screen ellipse

This second step is nonlinear because of perspective, so we use the Jacobian of the full camera-to-screen map.

So:

\[
\Sigma_c \to \Sigma_{\text{screen}}
\]

using

\[
\Sigma_{\text{screen}} \approx J_h \Sigma_c J_h^T
\]

\section{Full Chain from World Covariance}

If the world-to-camera linear part is $(W)$, then:

\[
\Sigma_c = W \Sigma_w W^T
\]

and then

\[
\Sigma_{\text{screen}} \approx J_h \Sigma_c J_h^T
\]

So:

\[
\Sigma_{\text{screen}} \approx J_h W \Sigma_w W^T J_h^T
\]

That is the full chain.

\section{Where Did the Skipped Steps Go?}

They are all inside $(h)$.

That is what I was compressing earlier.

Instead of spelling out:

\begin{itemize}
\item projection matrix
\item clip space
\item divide by $(w)$
\item NDC
\item viewport
\end{itemize}

I wrapped them into one function:

\[
h(x_c)
\]

Then $(J_h)$ is the Jacobian of that full chain.

So none of the steps disappeared. They were just bundled together.

\section{Two Equivalent Ways to Think About It}

\subsection{Way A: Direct Computer Vision Style}

Use:

\[
u = f_x \frac{X_c}{Z_c} + c_x,\qquad
v = f_y \frac{Y_c}{Z_c} + c_y
\]

Then compute the Jacobian directly from that.

\subsection{Way B: Full Graphics Pipeline Style}

Use:

\begin{enumerate}
\item projection matrix
\item clip space
\item perspective divide
\item NDC
\item viewport transform
\end{enumerate}

Then compute the Jacobian of the whole chain.

Both give the same idea.

\section{The Exact Process, No Skipped Steps}

Here is the full process as a chain.

\textbf{Point Pipeline}

\[
x_c \in \mathbb{R}^3
\]

\textbf{Step 1:} homogeneous lift

\[
x_c \to \tilde{x}_c =
\begin{bmatrix}
X_c\\
Y_c\\
Z_c\\
1
\end{bmatrix}
\]

\textbf{Step 2:} projection matrix

\[
\tilde{x}_{\text{clip}} = P \tilde{x}_c
\]

\textbf{Step 3:} perspective divide

\[
x_{\text{ndc}} =
\begin{bmatrix}
x_{\text{clip}} / w_{\text{clip}}\\
y_{\text{clip}} / w_{\text{clip}}\\
z_{\text{clip}} / w_{\text{clip}}
\end{bmatrix}
\]

\textbf{Step 4:} viewport transform

\[
s =
\begin{bmatrix}
u\\
v
\end{bmatrix}
\]

So:

\[
s = h(x_c)
\]

\textbf{Covariance Pipeline}

Start with camera-space covariance:

\[
\Sigma_c
\]

Take small 3D offsets $(\delta x)$ around the mean.

The exact screen offset is:

\[
\delta s = h(\mu_c + \delta x) - h(\mu_c)
\]

Approximate it linearly:

\[
\delta s \approx J_h \delta x
\]

Then covariance becomes:

\[
\Sigma_{\text{screen}} \approx J_h \Sigma_c J_h^T
\]

If you include world-to-camera:

\[
\Sigma_c = W \Sigma_w W^T
\]

then:

\[
\Sigma_{\text{screen}} \approx J_h W \Sigma_w W^T J_h^T
\]

\section{What the Code Is Conceptually Doing}

Yes, the code is conceptually doing all these steps.

But in implementation, it may:

\begin{itemize}
\item combine several steps
\item use a direct closed-form Jacobian
\item avoid building every intermediate quantity explicitly
\end{itemize}

So the math is the full story, even if the code is more compact.

\section{The Single-Sentence Summary}

Nothing was skipped mathematically: the steps were bundled into one full camera-to-screen function $(h)$, and the Jacobian $(J_h)$ is the derivative of that full chain, used to turn the 3D Gaussian covariance into a 2D screen-space covariance.

If you want, next I can do the \textbf{same full chain with actual numbers at every single stage}: camera point $\to$ clip $\to$ NDC $\to$ screen, and then compute $(J_h)$.

\end{document}
